\documentclass{article}

\usepackage{amsmath}
\usepackage{amsthm}
\usepackage{hyperref}

\newtheorem{definition}{Definition}
\newtheorem{lemma}{Lemma}
\newtheorem{theorem}{Theorem}
\newtheorem{corollary}{Corollary}

\title{A Hamming(7,4) Single-Error-Correction Target}
\author{EPFLemma workflow fixture}
\date{}

\begin{document}
\maketitle

\section{Motivation}

Hamming codes are a standard error-correcting mechanism from digital
communications and memory systems.  This fixture extracts a small but real
formalization target from the distance argument behind the binary Hamming
$(7,4)$ code introduced by Hamming~\cite{hamming1950}.  The target is compact:
it should fit in one generated Lean file, but it still requires the planner to
introduce definitions, choose a representation for binary words, split a
minimum-distance argument into lemmas, and connect that argument to uniqueness
of one-error correction.

The Lean project contains local setup definitions in
\texttt{DocFormalizationDemo.Setup}: binary bits as $\mathbb{F}_2$, words of
fixed length, Hamming weight, Hamming distance, the systematic Hamming $(7,4)$
encoder, a syndrome map, and a single-bit flip operation.

\section{Definitions}

\begin{definition}[Binary word]\label{def:binary_word}
A binary word of length $n$ is a function from the positions
$\{0,\ldots,n-1\}$ to the field $\mathbb{F}_2$.
\end{definition}

\begin{definition}[Hamming distance]\label{def:hamming_distance}
The Hamming distance $d(x,y)$ between two binary words of the same length is
the number of positions at which they differ.
\end{definition}

\begin{definition}[Systematic Hamming $(7,4)$ encoder]\label{def:hamming74_encoder}
For a message $(d_1,d_2,d_3,d_4) \in \mathbb{F}_2^4$, define the encoded word
$(p_1,p_2,d_1,p_3,d_2,d_3,d_4) \in \mathbb{F}_2^7$, where
\[
  p_1 = d_1+d_2+d_4,\qquad
  p_2 = d_1+d_3+d_4,\qquad
  p_3 = d_2+d_3+d_4.
\]
All additions are in $\mathbb{F}_2$.
\end{definition}

\section{Distance Argument}

\begin{lemma}[Zero syndrome for encoded words]\label{lem:encoded_syndrome_zero}
Every word produced by the encoder in Definition~\ref{def:hamming74_encoder}
has zero syndrome for the three parity checks
\[
  x_1+x_3+x_5+x_7,\qquad
  x_2+x_3+x_6+x_7,\qquad
  x_4+x_5+x_6+x_7.
\]
\end{lemma}

\begin{lemma}[No nonzero encoded word has weight one or two]\label{lem:no_small_nonzero_codeword}
If an encoded Hamming $(7,4)$ word is not the zero word, then its Hamming
weight is at least $3$.
\end{lemma}

\begin{theorem}[Minimum distance of the Hamming $(7,4)$ code]\label{thm:hamming74_min_distance}
If two encoded Hamming $(7,4)$ words come from distinct messages, then their
Hamming distance is at least $3$.
\end{theorem}

\section{Correction Property}

\begin{lemma}[Two radius-one balls around codewords are disjoint]\label{lem:radius_one_balls_disjoint}
Let $c_1$ and $c_2$ be encoded Hamming $(7,4)$ words.  If some received word
$r$ satisfies $d(r,c_1)\le 1$ and $d(r,c_2)\le 1$, then $c_1=c_2$.
\end{lemma}

\begin{theorem}[Single-error correction is unique]\label{thm:single_error_correction_unique}
For every received word produced from an encoded Hamming $(7,4)$ word by
flipping at most one bit, there is a unique encoded Hamming $(7,4)$ word within
Hamming distance at most $1$ of the received word.
\end{theorem}

\begin{corollary}[The transmitted message is unique]\label{cor:message_unique}
Under the same one-bit-error assumption, the original four-bit message is
uniquely determined by the received seven-bit word.
\end{corollary}

\begin{thebibliography}{9}
\bibitem{hamming1950}
R. W. Hamming,
\emph{Error Detecting and Error Correcting Codes},
Bell System Technical Journal 29(2), 147--160, 1950.
DOI: \url{https://doi.org/10.1002/j.1538-7305.1950.tb00463.x}.
\end{thebibliography}

\end{document}
